\documentclass[11pt,parskip=half]{scrartcl}

\usepackage[T1]{fontenc}
\usepackage{../common}
\usepackage{./plantuml}


\usepackage{csquotes} % To put text into “fancy quotes” with \textquote{bla bla}
\usepackage{epigraph} % For fancy wize quotes at the beginning of a section
\usepackage{eso-pic} % To draw a picture on top of a page
\usepackage{floatbytocbasic}

\usepackage{minted}
\usemintedstyle{colorful}

\usepackage{multicol}
\usepackage{pdfpages}
\usepackage{svg}

\renewcommand{\today}{\DTMdate{2026-04-29}}

\title{Trying out PlantUML diagrams in \LaTeX}

\author{Richard Pavy}
\date{\today}

\begin{document}

\maketitle

\tableofcontents
\listoffigures
\listoflistings

\section{Getting started with \texttt{PlantUML}}

This section contains sample \href{https://plantuml.com}{PlantUML} diagrams.

\subsection{Set Up}

Diagrams can be easily included by

\begin{minted}[linenos]{latex}
\usepackage{./plantuml}
...
% Includes diagram my-sequence-diagram.puml as a figure
% captioned "Sample UML diagram"
\includeplantuml{my-sequence-diagram}{Sample UML diagram}
\end{minted}

\subsection{Sample UML diagrams}

\begin{multicols}{2}
  \includeplantuml{participants}{Types of participants supported by PlantUML}
  \includeplantuml{arrows}{Types of arrows}
  \includeplantuml{groups}{Grouping message}
  \includeplantuml{anchors}{Anchors}
\end{multicols}

\subsection{A more complex example}

\includeplantuml{SSO}{Sample diagram of SSO handshake}

\subsection{A Wireframe example}

This example is configured with
\verb|\includeplantuml[width=\maxwidth]{puml}{Title}| to be as wide as possible
and fit the page.

\includeplantuml[width=\maxwidth]{salt}{S.A.L.T.}

\subsection{Inlined UML diagram}

UML diagrams can also be inlined in the \texttt{.tex} file, using:

\begin{minted}[linenos]{latex}
    \usepackage{./plantuml}
    ...
    % Includes inlined diagram as a figure captioned "Sample UML diagram"
    \begin{plantuml}{Alice and Bob}
        @startuml
        Alice -> Bob: Hello
        Bob -> Alice: Hi there
        @enduml
    \end{plantuml}
\end{minted}

\begin{plantuml}{Alice and Bob}
    @startuml
    !theme materia
    Alice -> Bob: Hello
    Bob -> Alice: Hi there
    @enduml
\end{plantuml}

\section{OIDC Authorization Code Flow}

\epigraph{Identity is no longer a place you log into, but a trust you carry across boundaries-quietly verified, never exposed.}{\textit{ChatGPT}}

\begin{plantuml}{OIDC Authorization Code Flow}
  @startuml
  !theme cerulean
  actor User
  participant "User Agent" as Browser
  participant "Client App" as Client
  participant "Authorization Server" as AS
  participant "Resource Server" as RS

  User -> Browser: Open protected page
  Browser -> Client: GET /protected
  Client -> Browser: Redirect to /authorize
  Browser -> AS: Authorization request\nclient_id, scope=openid, redirect_uri, state, nonce
  AS -> User: Authenticate and consent
  User -> AS: Credentials and consent
  AS -> Browser: Redirect with authorization code and state
  Browser -> Client: GET /callback?code=...&state=...
  Client -> AS: Exchange code for tokens
  AS -> Client: ID token, access token, refresh token
  Client -> AS: Validate ID token / fetch user info
  Client -> RS: API request with access token
  RS -> Client: Protected resource
  Client -> Browser: Authenticated session response
  @enduml
\end{plantuml}

\section{Source code snippets}

This is a \mint{html}|<h2>source code</h2>| snippet:

I can type \texttt{verbatim} text too.

\begin{listing}[H]
  \begin{minted}[linenos]{rust}
#[template]
#[html]
pub fn run() -> XElement {
    let demo = Demo::new();
    return div(
        key = "root",
        counter::counter_demo(demo.counter),
        timer::timer_demo(demo.enable_timer, demo.timer),
        api::api_demo(),
        attributes::attributes_demo(),
    );
}
\end{minted}
  \caption{Some Rust code}
  \label{listing:rust}
\end{listing}

\end{document}
